\section{Three-Curve Intersection and the Moment of Unified State}
\label{sec:curve-intersection}

% ═══════════════════════════════════════════════════════════════════
\subsection{Three Independent Charge-Redistribution Trajectories}
% ═══════════════════════════════════════════════════════════════════

A living system in a waking state maintains three separate, concurrent charge-redistribution trajectories, each operating on its own timescale and each decaying independently from its respective driving source.

\begin{definition}[Perception trajectory]
\label{def:perception-traj}
The perception trajectory $\mathcal{P}(t)$ is the charge-redistribution state of sensory processing circuits as they respond to current external input. It decays from peak sensory activation (when a stimulus arrives) toward baseline at a rate determined by sensory transduction kinetics and integration time constants.

Characteristic timescale: $\tau_{\text{perc}} \sim 100\text{--}500\,\text{ms}$ (visual integration window, auditory processing window).

The trajectory is driven by external sensory inputs: photons, pressure waves, chemical gradients, thermal energy.
\end{definition}

\begin{definition}[Thought trajectory]
\label{def:thought-traj}
The thought trajectory $\mathcal{T}(t)$ is the charge-redistribution state of deliberative circuits (prefrontal, parietal, premotor cortex) as they explore possible motor commands and predictions. It decays from peak exploration (when a decision is being made) toward a committed state at a rate determined by the slow-closure integration timescale.

Characteristic timescale: $\tau_{\text{thought}} \sim 1\text{--}5\,\text{s}$ (slow closure).

The trajectory is driven by internal goals, learned models, and the current attractor basin geometry.
\end{definition}

\begin{definition}[Memory trajectory]
\label{def:memory-traj}
The memory trajectory $\mathcal{M}(t)$ is the charge-redistribution state that encodes the history of where the system has been: the sequence of previous intersection points between perception and thought. It is the integral of past state-crossing events, accessible as proprioceptive return and temporal context in prefrontal and parietal circuits.

Characteristic timescale: $\tau_{\text{memory}} \sim 0.5\text{--}2\,\text{s}$ (working memory decay, contextual binding).

The trajectory is driven by the pattern of previous perception-thought pairings and by the system's trajectory through state space over the preceding seconds to minutes.
\end{definition}

\begin{remark}[Independence of the trajectories]
The three trajectories are not functions of one another. Perception can be present with thought absent (reflexive response to sudden threat). Thought can be present with perception absent (mental arithmetic with eyes closed). Memory can decay while both perception and thought remain active (focused attention with poor learning). They are topologically independent.
\end{remark}

% ═══════════════════════════════════════════════════════════════════
\subsection{Decay and Intersection}
% ═══════════════════════════════════════════════════════════════════

Each trajectory decays from its source toward a baseline or attractor state. The decay occurs on the trajectory's characteristic timescale.

\begin{definition}[Decay of a trajectory]
For trajectory $\mathcal{X}(t)$ driven by source $S(t)$, the decay is:
\begin{equation}
\frac{d\mathcal{X}}{dt} = -\frac{1}{\tau_X}(\mathcal{X} - \mathcal{X}_{\infty}) + \text{coupling terms}
\end{equation}
where $\tau_X$ is the decay timescale, $\mathcal{X}_{\infty}$ is the equilibrium or attractor, and coupling terms represent interaction with other trajectories.

Without external driving ($S = 0$), the trajectory asymptotically approaches $\mathcal{X}_{\infty}$.
\end{definition}

At any given moment during waking consciousness, all three trajectories are present in the charge-redistribution state space, each at a different level of decay. They intersect—that is, there exist isolated moments when the three trajectories align in a meaningful way.

\begin{definition}[Intersection point of three trajectories]
\label{def:intersection}
An intersection point $\mathcal{I}(t)$ is a moment at which:
\begin{enumerate}[label=(\roman*),leftmargin=*,itemsep=1pt]
\item The perception trajectory has decayed to a categorical state that can be labeled: ``I perceive X.''
\item The thought trajectory has decayed to a committed hypothesis that can be labeled: ``I think/expect Y.''
\item The memory trajectory provides a reference state such that the pairing (perception X, thought Y) is consistent with previous pairings.
\end{enumerate}

The intersection point is a region in charge-redistribution space, not a single point. But importantly, the intersection point is in NONE of the individual trajectory spaces. It is the point where the three curves meet.
\end{definition}

\begin{remark}[Intersection is not a state of any single system]
An intersection point cannot be described as ``a state of perception,'' ``a state of thought,'' or ``a state of memory'' alone. It is the alignment of all three. Remove any one trajectory, and the intersection point vanishes.
\end{remark}

% ═══════════════════════════════════════════════════════════════════
\subsection{Coherence Through Memory}
% ═══════════════════════════════════════════════════════════════════

At each moment, the three trajectories intersect at a point. At the next moment, they have each decayed further, and a new intersection point may form. For consciousness to be coherent—for the stream of intersection points to form a continuous, causally consistent narrative—the memory trajectory must validate each new intersection against the previous one.

\begin{theorem}[Intersection coherence via memory]
\label{thm:coherence}
Let $\mathcal{I}(t_n) = (\mathcal{P}(t_n), \mathcal{T}(t_n), \mathcal{M}(t_n))$ denote an intersection point at time $t_n$. Let $\mathcal{I}(t_{n+1})$ denote the next intersection point at time $t_{n+1} = t_n + \Delta t$, where $\Delta t \sim \tau_{\text{perc}}$ (perception decay timescale).

The sequence of intersection points is causally coherent if $\mathcal{M}(t_{n+1})$ encodes information about $\mathcal{I}(t_n)$. Specifically:
\begin{equation}
\mathcal{M}(t_{n+1}) = f(\mathcal{I}(t_n), \mathcal{M}(t_n)),
\end{equation}
where $f$ is a function that relates the previous intersection to the current memory state. This ensures that the current intersection $\mathcal{I}(t_{n+1})$ can be validated against the previous one.
\end{theorem}

\begin{proof}
Without the memory term, each intersection point would be independent. Perception at $t_{n+1}$ might be "cup," thought might be "weapon," and there would be no constraint on whether this pairing is reasonable. But with memory encoding the previous intersection (where "cup" was paired with "cylindrical object" and "drinkable from"), the current intersection can be validated: "Yes, seeing a cup-shaped object and thinking 'cup' is consistent with my previous experience."

If $\mathcal{M}$ did not carry forward information about $\mathcal{I}(t_n)$, each moment would feel disconnected, and learning would be impossible.
\end{proof}

\begin{corollary}[Trajectory history is reference]
The memory trajectory is the integrated history of previous intersection points. It is not a storage system but rather a reference. It answers the question: "Does this intersection (this perception-thought pairing) make sense given where I've been?"
\end{corollary}

\subsection{Absence of a Trajectory: What Happens}

The structure of three-curve intersection predicts specific deficits when any one trajectory is absent or disrupted.

\begin{proposition}[Missing perception trajectory]
If sensory input is gated off (eyes closed, proprioceptive return inhibited, as in REM sleep), the perception trajectory becomes flat or absent. The thought trajectory continues to decay, but without a corresponding perception trajectory, the thought has no external constraint. No meaningful intersection points form. The thought trajectory oscillates in high-dimensional space without convergence to a realistic attractor.
\end{proposition}

\begin{proposition}[Missing thought trajectory]
If deliberative circuits are offline (deep coma, under general anesthesia), the thought trajectory flattens. Even if perception is present (sensory input arrives), without thought to pair with it, perception alone cannot form meaningful intersection points. The system exhibits reflexive responses to stimuli without unified experience.
\end{proposition}

\begin{proposition}[Missing memory trajectory]
If proprioceptive feedback is severed (complete deafferentation) or if working memory is disrupted, the memory trajectory cannot carry forward information about previous intersections. Each intersection point becomes isolated. The system cannot validate that successive perceptions and thoughts are coherent. Without this validation, learning is impossible and behavior becomes disconnected.
\end{proposition}

\subsection{The Intersection Point is Consciousness}

\begin{principle}[Consciousness as intersection manifold]
\label{prin:consciousness-intersection}
Consciousness is not a state, property, or substance. It is the manifold of intersection points where the three charge-redistribution trajectories (perception, thought, memory) align in a causally coherent sequence.

At each moment of consciousness:
\begin{itemize}
\item A perception trajectory has decayed to a categorical state (``I perceive'').
\item A thought trajectory has decayed to a committed hypothesis (``I think'').
\item A memory trajectory provides continuity (``This makes sense given where I was'').
\item The three align at an intersection point.
\end{itemize}

The intersection point itself lies in none of the three trajectory spaces. It is the alignment itself—a topological event, not a state within any component system.

Successive intersection points form a stream. This stream is the subjective experience of consciousness.
\end{principle}

\begin{remark}[Why all three are necessary]
Remove perception → thought unconstrained, oscillates wildly (dreams, delirium).
Remove thought → perception uninterpreted, reflexive only (coma, deep anesthesia).
Remove memory → intersections isolated, no coherence (deafferentation, amnesia, severe dissociation).

All three must be present and aligned for consciousness to emerge.
\end{remark}

% ═══════════════════════════════════════════════════════════════════
\subsection{Charge Rate at the Intersection}
% ═══════════════════════════════════════════════════════════════════

The charge redistribution associated with an intersection point is the sum of charge rates from all three trajectories at the moment of intersection.

\begin{definition}[Intersection charge rate]
At an intersection point $\mathcal{I}(t)$:
\begin{equation}
Q_{\text{intersection}}(t) = Q_{\text{perc}}(t) + Q_{\text{thought}}(t) + Q_{\text{memory}}(t)
\end{equation}

where:
\begin{align}
Q_{\text{perc}}(t) &= \text{charge from sensory transduction and integration} \\
Q_{\text{thought}}(t) &= \text{charge from deliberative circuits} \approx 133.5\,\text{mC/s (from Section 6)} \\
Q_{\text{memory}}(t) &= \text{charge from working memory and contextual binding}
\end{align}
\end{definition}

The intersection charge rate is not static. As each trajectory decays or is renewed by new input, the intersection charge rate changes. The variation in intersection charge rate is the signature of the stream of consciousness.

\subsection{Continuity and Uniqueness}

\begin{proposition}[Poincaré deviation at intersections]
Just as no two motor-unit activations are identical (Theorem 5.6), no two intersection points are identical. Each moment of consciousness is unique. The three trajectories never return to exactly the same configuration.

This is why:
\begin{itemize}
\item You cannot have the exact same thought twice.
\item You cannot perceive an object identically twice.
\item Each conscious moment is a singular, unrepeatable event.
\item Memory must bridge the gap: ``I was here before, but not exactly like this.''
\end{itemize}
\end{proposition}

\begin{proposition}[Continuity of the intersection stream]
Although each intersection point is unique (Poincaré deviation), the sequence of intersection points forms a continuous trajectory through state space if and only if the memory trajectory successfully encodes information from each intersection to the next.

If memory is interrupted, the stream becomes fragmented. If memory is absent, no stream forms.
\end{proposition}

\subsection{Summary: The Foundation for Unified Experience}

The three-curve intersection model predicts:

\begin{enumerate}[leftmargin=1.3em]
\item \textbf{Consciousness requires all three:} Remove perception, thought, or memory, and consciousness ceases. No exceptions.

\item \textbf{Consciousness is dynamic:} It is not a state but a sequence of intersection points, each unique, each requiring the others.

\item \textbf{Memory is essential for coherence:} Without memory validating each new intersection against the previous one, there is no unified narrative—only fragmented moments.

\item \textbf{The intersection point is not in any trajectory:} Consciousness is not a property of perception alone, thought alone, or memory alone. It emerges from their alignment.

\item \textbf{Charge redistribution marks the intersection:} The combined charge from all three trajectories at their intersection point is the quantifiable correlate of the conscious moment.

\item \textbf{Each moment is unique:} Poincaré deviation ensures no two intersections are identical. The stream is a sequence of singular events.

\end{enumerate}

This foundation now permits the analysis of what happens when individual trajectories are isolated, modulated, or disrupted—which will be the focus of the remaining sections on perception, thought, memory, and their pathological dissociations.